\documentclass[11pt,a4paper]{article}

\usepackage[margin=1in]{geometry}
\usepackage{amsmath,amssymb}
\usepackage{booktabs}
\usepackage{enumitem}
\usepackage{titlesec}
\usepackage{hyperref}
\usepackage{xcolor}
\usepackage{fancyhdr}
\usepackage{parskip}
\usepackage{microtype}
\usepackage{tikz}

% Colors
\definecolor{embed}{RGB}{0,100,0}
\definecolor{intuition}{RGB}{0,50,130}
\definecolor{falstad}{RGB}{130,50,0}

% Section formatting
\titleformat{\section}{\Large\bfseries}{}{0em}{}[\vspace{-0.5em}\rule{\textwidth}{0.4pt}]
\titleformat{\subsection}{\large\bfseries}{}{0em}{}
\titleformat{\subsubsection}{\normalsize\bfseries}{}{0em}{}

% Header/footer
\pagestyle{fancy}
\fancyhf{}
\fancyhead[L]{\small\textit{Art of Electronics --- Ch.\ 1 Foundations}}
\fancyhead[R]{\small\thepage}
\renewcommand{\headrulewidth}{0.4pt}

% Custom environments
\newcommand{\intuition}[1]{\par\textcolor{intuition}{\textbf{Intuition:}} #1}
\newcommand{\embedded}[1]{\par\textcolor{embed}{\textbf{Embedded tie-in:}} \textit{#1}}

% Simple circuit drawing helpers (no circuitikz needed)
\newcommand{\resistor}[2]{% start pos, label
  % zigzag resistor drawn manually
}
\tikzset{
  gnd/.style={draw, fill=white, inner sep=0pt,
    path picture={
      \draw (path picture bounding box.north) -- (path picture bounding box.south);
      \draw ([yshift=-1pt]path picture bounding box.south west) -- ([yshift=-1pt]path picture bounding box.south east);
    }
  }
}

\title{\vspace{-2em}\textbf{Chapter 1 --- Foundations}\\[0.3em]
\large Companion Notes for \textit{The Art of Electronics}, 3rd ed.\\
Sections 1.1--1.10, pp.\ 1--67}
\author{}
\date{}

\begin{document}
\maketitle
\thispagestyle{fancy}

%=============================================================================
\section{Voltage, Current, Resistance (\S1.2)}
%=============================================================================

\subsection{Ohm's Law and Kirchhoff's Laws}

The three rules that govern all resistive circuits:

\begin{align}
V &= IR \tag{Ohm's law} \\
\sum I_{\text{in}} &= \sum I_{\text{out}} \tag{KCL} \\
\sum V_{\text{loop}} &= 0 \tag{KVL}
\end{align}

\intuition{Voltage is the \emph{cause} (pressure), current is the \emph{effect} (flow), resistance is the \emph{opposition}. Current flows \emph{through} things; voltage appears \emph{across} things. Never say ``voltage through a resistor''---that's nonsense.}

\subsection{Resistors: Series and Parallel}

\textbf{Series:}
\[R_{\text{total}} = R_1 + R_2 + \cdots\]

\textbf{Parallel:}
\[R_{\text{total}} = \frac{R_1 R_2}{R_1 + R_2} \quad\text{(two resistors)}, \qquad
\frac{1}{R_{\text{total}}} = \frac{1}{R_1} + \frac{1}{R_2} + \cdots \quad\text{(general)}\]

\textbf{Shortcut \#1:} A large resistor in series (parallel) with a small one has roughly the resistance of the larger (smaller) one. Series to trim up, parallel to trim down.

\textbf{Shortcut \#2:} Think in equal-resistor multiples. $5\text{k} \| 10\text{k}$? That's three $10\text{k}$'s in parallel $= 10\text{k}/3 = 3.33\text{k}$. In general, $n$ equal resistors in parallel $= R/n$.

\subsection{Power Dissipation}

\[
P = IV = I^2 R = \frac{V^2}{R}
\]

All three forms are Ohm's law substitutions of $P = IV$. Pick whichever matches what you know. A $\frac{1}{4}\,$W resistor above $1\,\text{k}\Omega$ can never blow from a 15\,V battery (see Exercise~1.5---prove it to yourself).

\subsection{Voltage Dividers}

The most important circuit fragment in electronics (Figure~1.6, p.~7). Two resistors in series, output from the middle:

\begin{equation}
\boxed{V_{\text{out}} = V_{\text{in}} \cdot \frac{R_2}{R_1 + R_2}}
\end{equation}

\intuition{The output is the fraction of total resistance ``below'' the tap. If $R_2$ is $\frac{1}{3}$ of the total, you get $\frac{1}{3}$ of the input. Any node in a circuit can be modeled as a divider between the driving impedance above and the load impedance below.}

\embedded{Pull-up/pull-down resistors on MCU GPIO pins are one leg of a voltage divider. A $10\,\text{k}\Omega$ pull-up to 3.3\,V with a button to ground switches between 3.3\,V (open) and 0\,V (pressed). ADC input scaling (e.g., measuring 12\,V with a 3.3\,V ADC) uses a resistive divider.}

\subsection{Voltage and Current Sources}

\begin{itemize}[nosep]
\item \textbf{Voltage source:} Maintains $V$ across terminals. $R_{\text{int}}$ ideally zero. A 9\,V battery $\approx$ perfect 9\,V source $+$ $3\,\Omega$ series resistance.
\item \textbf{Current source:} Maintains $I$ through terminals. $R_{\text{int}}$ ideally infinite. ``Likes'' short circuits, ``hates'' open circuits.
\end{itemize}

\subsection{Th\'{e}venin Equivalent Circuit}

Any two-terminal network of resistors and voltage sources is equivalent to a single voltage source $V_{\text{Th}}$ in series with a single resistor $R_{\text{Th}}$ (Figures~1.11--1.13, pp.~9--11):

\begin{equation}
V_{\text{Th}} = V_{\text{open\,circuit}}, \qquad
R_{\text{Th}} = \frac{V_{\text{open\,circuit}}}{I_{\text{short\,circuit}}}
\end{equation}

For a voltage divider:
\begin{equation}
V_{\text{Th}} = V_{\text{in}} \cdot \frac{R_2}{R_1 + R_2}, \qquad
R_{\text{Th}} = R_1 \| R_2 = \frac{R_1 R_2}{R_1 + R_2}
\end{equation}

\intuition{Th\'{e}venin tells you how ``stiff'' a source is. Low $R_{\text{Th}}$ = output barely sags under load. To avoid loading, keep $R_{\text{load}} \gg R_{\text{Th}}$ (rule of thumb: $10\times$ or more). See Figure~1.14 (p.~11) for the attenuation vs.\ load-ratio graph.}

\subsection{Small-Signal (Dynamic) Resistance}

For nonlinear devices (diodes, zeners), the ratio $V/I$ isn't constant. Instead, we use the slope of the $V$--$I$ curve:

\begin{equation}
r_{\text{dyn}} = \frac{dV}{dI} = \frac{\Delta V}{\Delta I}
\end{equation}

\intuition{A zener at its rated current might have $V_Z = 5.1\,$V but $r_{\text{dyn}} = 10\,\Omega$. A 10\% change in current ($\Delta I = 1\,$mA around 10\,mA) causes $\Delta V = 10 \times 0.001 = 10\,$mV---only 0.2\% regulation. For small signals, the zener \emph{acts like} a $10\,\Omega$ resistor in a divider with the series dropping resistor. This concept is essential for transistor circuits in Ch.~2.}

\subsection{Circuit Loading}

When a load is attached, the output drops because current flows through the source's Th\'{e}venin resistance. Always ensure $R_{\text{load}} \gg R_{\text{source}}$.

\embedded{The STM32 ADC has $\sim 6\,\text{k}\Omega$ input impedance during sampling---your source must drive that without sagging. Op-amp outputs have $R_{\text{Th}}$ of milliohms. A GPIO pin driving HIGH has $R_{\text{Th}} \approx 25$--$50\,\Omega$.}

%=============================================================================
\section{Signals (\S1.3)}
%=============================================================================

\subsection{Sinusoidal Signals}

\begin{equation}
V(t) = A \sin(2\pi f t + \phi) = A \sin(\omega t + \phi), \qquad \omega = 2\pi f
\end{equation}

Key relationships:
\begin{align}
V_{\text{rms}} &= \frac{V_{\text{peak}}}{\sqrt{2}} = 0.707 \, V_{\text{peak}} \quad \text{(sinewaves only)} \\
V_{\text{pp}} &= 2 \, V_{\text{peak}}
\end{align}

US mains: $120\,$V rms, $170\,$V peak, $60\,$Hz.

\intuition{Sinewaves are special because linear circuits preserve their shape---a sinewave in always produces a sinewave out (possibly with different amplitude and phase). This is why frequency-domain analysis works.}

\subsection{Decibels}

\begin{equation}
\text{dB} = 10 \log_{10}\!\frac{P_2}{P_1} = 20 \log_{10}\!\frac{A_2}{A_1}
\end{equation}

\textbf{Landmarks to memorize:}

\begin{center}
\begin{tabular}{@{}rcc@{}}
\toprule
dB & Voltage ratio & Power ratio \\
\midrule
0  & 1        & 1    \\
3  & $\sqrt{2} \approx 1.41$ & 2    \\
6  & 2        & 4    \\
10 & 3.16     & 10   \\
20 & 10       & 100  \\
$-3$  & 0.707    & 0.5  \\
$-20$ & 0.1      & 0.01 \\
\bottomrule
\end{tabular}
\end{center}

Everything else is addition: $+26\,$dB $= +20 + 6 = 10\times \cdot 2\times = 20\times$ voltage.

\subsection{Other Signal Types}

\begin{itemize}[nosep]
\item \textbf{Square wave:} Digital signals. $V_{\text{peak}} = V_{\text{rms}}$. Characterized by rise time (10\% to 90\%).
\item \textbf{Triangle/ramp:} Basis of sweep circuits, sawtooth oscillators, ADC conversion.
\item \textbf{Noise:} Random thermal fluctuations (Johnson noise). Specified as power spectral density or rms voltage.
\item \textbf{Pulses:} Defined by amplitude, width, repetition rate, duty cycle.
\end{itemize}

\subsection{Logic Levels}

\begin{center}
\begin{tabular}{@{}lcc@{}}
\toprule
Parameter & 3.3\,V LVCMOS & 5\,V TTL \\
\midrule
$V_{\text{OH}}$ (output high) & $\sim$3.3\,V & $\sim$5\,V \\
$V_{\text{OL}}$ (output low)  & $\sim$0\,V   & $\sim$0\,V \\
$V_{\text{IH}}$ (input high threshold)  & 1.5\,V & 2.0\,V \\
$V_{\text{IL}}$ (input low threshold)   & 0.8\,V & 0.8\,V \\
\bottomrule
\end{tabular}
\end{center}

\embedded{PWM is a rectangular signal with variable duty cycle---the ``signal'' is the average voltage ($D \cdot V_{\text{high}}$). UART idle is HIGH; start bit pulls LOW. When interfacing 3.3\,V and 5\,V logic, use level shifters or series resistors with clamp diodes.}

%=============================================================================
\section{Capacitors \& RC Circuits (\S1.4)}
%=============================================================================

\subsection{Capacitor Fundamentals}

\begin{align}
Q &= CV \tag{charge stored} \\
I &= C\,\frac{dV}{dt} \tag{current $\propto$ rate of voltage change} \\
U &= \tfrac{1}{2}\,CV^2 \tag{stored energy}
\end{align}

\intuition{A capacitor resists changes in voltage. At DC it's an open circuit (no $dV/dt$ = no current). At high frequencies it's nearly a short (rapid voltage changes = large current).}

\textbf{Combining capacitors} (opposite of resistors!):
\[
\text{Parallel: } C_{\text{total}} = C_1 + C_2 + \cdots \qquad
\text{Series: } \frac{1}{C_{\text{total}}} = \frac{1}{C_1} + \frac{1}{C_2} + \cdots
\]

\subsection{Bypass vs.\ Coupling --- The Two Most Common Uses}

\begin{itemize}[nosep]
\item \textbf{Bypass (decoupling):} Cap from power pin to ground. Local charge reservoir for transient demands. Short circuit to high-frequency noise.
\item \textbf{Coupling (blocking):} Cap in series with signal path. Passes AC, blocks DC bias. Acts as highpass filter.
\end{itemize}

\embedded{Every MCU datasheet specifies bypass caps---typically $100\,$nF ceramic on each $V_{\text{DD}}$ pin, plus $10\,\mu$F bulk nearby. Without them, switching currents cause ground bounce and supply droop that corrupt logic.}

\subsection{RC Time Constant and Charge/Discharge}

\begin{equation}
\boxed{\tau = RC}
\end{equation}

\begin{align}
\text{Discharge:} \quad V(t) &= V_0 \, e^{-t/RC} \\
\text{Charge:} \quad V(t) &= V_f \left(1 - e^{-t/RC}\right)
\end{align}

See Figures~1.32--1.34 (pp.~21--22).

\textbf{The $5\tau$ rule:}

\begin{center}
\begin{tabular}{@{}cc@{}}
\toprule
Time & \% of final value (charging) \\
\midrule
$1\,RC$ & 63\% \\
$2\,RC$ & 86\% \\
$3\,RC$ & 95\% \\
$5\,RC$ & 99.3\% \\
\bottomrule
\end{tabular}
\end{center}

\intuition{After one time constant, the cap is 63\% charged. After five, it's done. A $1\,\mu$F cap across $1\,\text{k}\Omega$ has $\tau = 1\,$ms; fully settled in $\sim$5\,ms.}

\subsection{Differentiator (Edge Detector)}

Cap in series with input, resistor to ground (Figure~1.41B, p.~25):

\begin{equation}
V_{\text{out}} \approx RC \, \frac{dV_{\text{in}}}{dt} \qquad \text{(when } RC \ll T_{\text{signal}}\text{)}
\end{equation}

\intuition{Feed in a square wave, get spikes at the edges. The cap passes fast changes and blocks slow ones. This is a highpass filter in the frequency domain, a differentiator in the time domain---same circuit, two viewpoints.}

\subsection{Integrator (Averager)}

Resistor in series with input, cap to ground (Figure~1.46B, p.~27):

\begin{equation}
V_{\text{out}} \approx \frac{1}{RC} \int V_{\text{in}} \, dt \qquad \text{(when } V_{\text{out}} \ll V_{\text{in}}\text{)}
\end{equation}

\intuition{Feed in a square wave, get a triangle (ramp up, ramp down). The cap smooths out rapid changes, keeping only the slow trend. This is the lowpass filter in the frequency domain.}

%=============================================================================
\section{Inductors \& Transformers (\S1.5)}
%=============================================================================

\subsection{Inductor Fundamentals}

\begin{align}
V &= L\,\frac{dI}{dt} \tag{voltage $\propto$ rate of current change} \\
U &= \tfrac{1}{2}\,LI^2 \tag{stored energy in magnetic field}
\end{align}

\intuition{An inductor resists changes in current---the dual of a capacitor. If you suddenly open a switch carrying inductor current, it generates a huge voltage spike (``inductive kick'') to force the current to keep flowing.}

\subsection{Volt-Second Balance}

In steady-state periodic operation, the average voltage across an inductor must be zero:
\[
\langle V_L \rangle = \frac{1}{T}\int_0^T V_L\,dt = 0
\]
Otherwise the current would ramp without limit. This constraint governs all switching power converters.

\subsection{Transformers}

AC on the primary induces AC on the secondary, scaled by the turns ratio (Figure~1.53, p.~30):

\begin{equation}
\frac{V_{\text{sec}}}{V_{\text{pri}}} = \frac{N_{\text{sec}}}{N_{\text{pri}}} = n, \qquad
\frac{I_{\text{sec}}}{I_{\text{pri}}} = \frac{1}{n}, \qquad
Z_{\text{reflected}} = n^2 Z_{\text{load}}
\end{equation}

\intuition{Transformers trade voltage for current while conserving power, and provide galvanic isolation (no DC path between primary and secondary).}

\subsection{Buck \& Boost Preview}

\begin{itemize}[nosep]
\item \textbf{Buck (step-down):} Switch alternates inductor between $V_{\text{in}}$ and ground. $V_{\text{out}} = D \cdot V_{\text{in}}$, where $D$ is the duty cycle (Figure~1.52A).
\item \textbf{Boost (step-up):} Inductor stores energy from $V_{\text{in}}$, switch forces it through diode to output. $V_{\text{out}} = V_{\text{in}} / (1-D)$ (Figure~1.52B).
\end{itemize}

\embedded{Every modern MCU board has a switching regulator (usually buck) to convert USB 5\,V to 3.3\,V core. Ferrite beads on digital lines are tiny inductors---high impedance at RF, transparent at DC. Always include a flyback diode across relay/motor coils---without it the inductive kick destroys your MOSFET.}

%=============================================================================
\section{Diodes \& Diode Circuits (\S1.6)}
%=============================================================================

\subsection{Diode $V$--$I$ Curve}

A diode is a two-terminal passive nonlinear device (Figure~1.55, p.~31; Table~1.1, p.~32):

\begin{itemize}[nosep]
\item \textbf{Forward biased:} Current flows when $V_{\text{anode}} - V_{\text{cathode}} \gtrsim 0.6\,$V (Si) or $\sim$0.3\,V (Schottky).
\item \textbf{Reverse biased:} Essentially no current (nA) until breakdown (50--1000\,V).
\item A diode is \emph{not} a resistor. It doesn't obey Ohm's law. If you put diodes in a circuit, it won't have a Th\'{e}venin equivalent.
\end{itemize}

\subsection{Rectification}

Converting AC to DC using diodes (Figures~1.56--1.60, pp.~32--33):

\begin{itemize}[nosep]
\item \textbf{Half-wave:} Single diode, passes only positive half-cycles.
\item \textbf{Full-wave bridge:} Four diodes, passes both half-cycles. The standard for power supplies.
\end{itemize}

\textbf{Ripple voltage} (the key design equation):
\begin{equation}
\Delta V = \frac{I_{\text{load}}}{fC} \quad \text{(half-wave)}, \qquad
\Delta V = \frac{I_{\text{load}}}{2fC} \quad \text{(full-wave)}
\end{equation}

\intuition{Don't memorize these---derive from $I = C\,dV/dt$ with $\Delta t \approx 1/f$. The cap charges to peak during diode conduction, then discharges into the load between peaks. Bigger cap = less ripple but higher peak charging currents.}

\subsection{Zener Diodes and Voltage Regulation}

A zener operates in reverse breakdown at a specified voltage (Figure~1.16, p.~12). For \emph{changes} in voltage, the circuit behaves as a voltage divider with $r_{\text{dyn}}$:

\begin{equation}
\Delta V_{\text{out}} = \frac{r_{\text{dyn}}}{R + r_{\text{dyn}}} \, \Delta V_{\text{in}}
\end{equation}

\intuition{Low-voltage zeners ($<3.3\,$V) perform poorly (Figure~1.17, p.~13)---use IC voltage references instead.}

\subsection{Diode Circuit Applications}

\begin{itemize}[nosep]
\item \textbf{Clamps:} Limit signal voltage to a fixed range (Figure~1.72, p.~36).
\item \textbf{Limiters:} Back-to-back diodes clip to $\pm 0.6\,$V. Input protection (Figure~1.78, p.~37).
\item \textbf{OR gates (diode logic):} Pass the higher of two voltages. Battery backup (Figure~1.71, p.~36).
\item \textbf{Flyback diode:} Across inductive loads. \emph{Always.} 1N4004 works for nearly all cases (Figure~1.84, p.~39).
\end{itemize}

\embedded{Every CMOS GPIO has ESD protection diodes to $V_{\text{DD}}$ and GND. If your signal exceeds $V_{\text{DD}} + 0.6\,$V or goes below GND $- 0.6\,$V, current flows through these---can latch up and destroy the chip. Schottky diodes in power paths prevent back-feeding. Reverse polarity protection: series Schottky (simple, wastes 0.3\,V) or P-MOSFET (better, near-zero drop).}

%=============================================================================
\section{Impedance \& Reactance (\S1.7)}
%=============================================================================

\subsection{Reactance}

The frequency-dependent ``resistance'' of reactive components:

\begin{equation}
\boxed{X_C = \frac{1}{\omega C} = \frac{1}{2\pi f C}} \qquad \text{(decreases with }f\text{)}
\end{equation}
\begin{equation}
\boxed{X_L = \omega L = 2\pi f L} \qquad \text{(increases with }f\text{)}
\end{equation}

\intuition{Capacitors are ``short'' at high frequencies, ``open'' at DC. Inductors are the opposite. This is why caps bypass high-frequency noise and inductors choke it.}

\subsection{Complex Impedance}

\begin{equation}
\mathbf{Z}_R = R, \qquad
\mathbf{Z}_C = \frac{1}{j\omega C} = -\frac{j}{\omega C}, \qquad
\mathbf{Z}_L = j\omega L
\end{equation}

\textbf{Generalized Ohm's law:}
\begin{equation}
\mathbf{V} = \mathbf{I}\,\mathbf{Z}
\end{equation}

Series and parallel combination rules work identically to resistors, using complex arithmetic. The magnitude $|\mathbf{Z}|$ gives the voltage-to-current amplitude ratio; the angle gives the phase shift.

\subsection{RC Lowpass Filter}

$R$ in series, $C$ to ground (Figure~1.90, p.~42; Bode plot Figure~1.104, p.~50):

\begin{equation}
\frac{V_{\text{out}}}{V_{\text{in}}} = \frac{1}{\sqrt{1 + (2\pi f R C)^2}}, \qquad
\boxed{f_{3\text{dB}} = \frac{1}{2\pi RC}}
\end{equation}

Above $f_{3\text{dB}}$: rolls off at $-20\,$dB/decade ($-6\,$dB/octave). Phase: $0^\circ \to -45^\circ$ at $f_{3\text{dB}} \to -90^\circ$.

\subsection{RC Highpass Filter}

$C$ in series, $R$ to ground (Figure~1.92, p.~43):

\begin{equation}
\frac{V_{\text{out}}}{V_{\text{in}}} = \frac{2\pi f R C}{\sqrt{1 + (2\pi f R C)^2}}, \qquad
f_{3\text{dB}} = \frac{1}{2\pi RC}
\end{equation}

Same breakpoint formula. This is the ``blocking capacitor'' configuration.

\subsection{LC Resonance}

When an inductor and capacitor are combined, they exchange energy at a natural frequency (Figures~1.107--1.110, pp.~53--54):

\begin{equation}
\boxed{f_0 = \frac{1}{2\pi\sqrt{LC}}}
\end{equation}

\begin{itemize}[nosep]
\item \textbf{Parallel LC (``tank''):} Maximum impedance at $f_0$. Peak in frequency response.
\item \textbf{Series LC (``trap''):} Minimum impedance at $f_0$. Notch in frequency response.
\item \textbf{Quality factor:} $Q = f_0 / \Delta f_{3\text{dB}}$. Higher $Q$ = sharper peak = more selective.
\end{itemize}

\subsection{Power Factor}

\begin{equation}
\text{power factor} = \cos\phi = \frac{P_{\text{real}}}{V_{\text{rms}} \cdot I_{\text{rms}}}
\end{equation}

A purely reactive circuit ($\phi = 90^\circ$) has power factor = 0---current flows but no power is consumed.

\embedded{Anti-aliasing filter before an ADC: set $f_{3\text{dB}} \leq f_{\text{sample}}/2$. For a 1\,MSPS ADC, use e.g.\ $R = 1\,\text{k}\Omega$, $C = 1.6\,$nF ($f_{3\text{dB}} \approx 100\,$kHz). LC filters in switching supplies have much sharper cutoffs than RC. A $100\,$nF ceramic bypass cap has $|\mathbf{Z}| = 1.6\,\Omega$ at 1\,MHz but $160\,\Omega$ at 10\,kHz.}

%=============================================================================
\section{AM Radio Capstone (\S1.8)}
%=============================================================================

A circuit that ties together every concept from Chapter~1 (Figure~1.114, p.~56):

\begin{enumerate}[nosep]
\item \textbf{Antenna} picks up all AM stations plus noise.
\item \textbf{Parallel LC tank} (variable $C_1$ + inductor $L$) tuned to desired station (520--1720\,kHz). High $Q$ amplifies the selected station and rejects others.
\item \textbf{Diode $D$} half-wave rectifies the modulated carrier.
\item \textbf{$R_1$} provides a DC load so the rectified output follows the envelope.
\item \textbf{$C_2$} smooths the RF carrier. Time constant $R_1 C_2$ is long vs.\ carrier period ($\sim 1\,\mu$s) but short vs.\ audio period ($\sim 200\,\mu$s).
\item \textbf{Audio amplifier} drives the speaker.
\end{enumerate}

\intuition{The diode + RC filter is an ``envelope detector''---it extracts slow amplitude variations while discarding the fast carrier, the same principle as power-supply filtering but at RF frequencies. The LC tank not only selects but \emph{amplifies}---voltage across the tank at resonance is $Q$ times the antenna voltage (Figures~1.115--1.116, p.~56).}

%=============================================================================
\section{Other Passive Components (\S\S1.9--1.10)}
%=============================================================================

\subsection{Switches}

\begin{itemize}[nosep]
\item \textbf{SPST:} Simple on/off (Figure~1.119, p.~58).
\item \textbf{SPDT:} Selects between two connections. Three-way light switches.
\item \textbf{DPDT:} Two SPDT switches ganged. Can reverse motor polarity.
\item \textbf{Momentary (pushbutton):} NO or NC. Returns to default when released (Figure~1.120).
\item \textbf{Rotary:} Multiple positions. Shorting (make-before-break) or non-shorting.
\end{itemize}

\subsection{Relays}

Electrically controlled switches. A coil energizes an electromagnet that pulls contacts closed (\S1.9.2, p.~59):

\begin{itemize}[nosep]
\item \textbf{Electromechanical:} Complete galvanic isolation. Coil voltages 3--115\,V.
\item \textbf{Solid-state relay (SSR):} LED + phototransistor/triac. No bounce, longer life, faster.
\end{itemize}

\embedded{Never drive a relay coil directly from an MCU GPIO---too much current (10--100\,mA) and the inductive kick damages the pin. Use an NPN or N-MOSFET as a switch, with a flyback diode across the coil.}

\subsection{Connectors}

\begin{center}
\begin{tabular}{@{}ll@{}}
\toprule
Type & Use \\
\midrule
BNC & Lab instruments, 50/75\,$\Omega$ coaxial \\
SMA/SMB & RF boards, miniature coaxial \\
D-sub (DB-9, DB-25) & Serial ports, legacy I/O \\
Pin headers (0.1$''$/2.54\,mm) & Dev boards, prototyping \\
USB (A, B, Micro-B, Type-C) & Power and data \\
Edge connectors & PCIe, memory DIMMs \\
\bottomrule
\end{tabular}
\end{center}

\subsection{LEDs and Indicators}

LEDs are diodes with forward drop 1.5--2\,V (red/green) or 3--3.6\,V (blue/white). Typical $I$: 2--10\,mA.

\textbf{Series current-limiting resistor:}
\begin{equation}
R = \frac{V_{\text{supply}} - V_{\text{LED}}}{I_{\text{LED}}}
\qquad \text{Example: } R = \frac{3.3 - 2.0}{0.005} = 260\,\Omega \to 270\,\Omega
\end{equation}

\subsection{Variable Components}

\begin{itemize}[nosep]
\item \textbf{Potentiometers:} Variable resistors with a wiper. Don't use as precision resistors---poor stability. Use fixed $+$ small trim pot.
\item \textbf{Varactors:} Voltage-variable capacitors (reverse-biased diodes). Used in PLLs, VCOs, automatic frequency tuning.
\end{itemize}

\subsection{SMT Packages}

\begin{center}
\begin{tabular}{@{}llll@{}}
\toprule
Imperial & Metric & Size (mm) & Notes \\
\midrule
0402 & 1005 & $1.0 \times 0.5$ & Tweezers required \\
0603 & 1608 & $1.6 \times 0.8$ & Common for production \\
0805 & 2012 & $2.0 \times 1.25$ & Good for hand soldering \\
1206 & 3216 & $3.2 \times 1.6$ & Easy to hand solder \\
\bottomrule
\end{tabular}
\end{center}

%=============================================================================
\section{Falstad Simulation Suggestions}
%=============================================================================

\begin{enumerate}[nosep]
\item \textbf{Voltage divider with loading:} $10\text{k}/10\text{k}$ from 10\,V. Measure unloaded (5\,V), then add 10k load (3.33\,V). Verify Th\'{e}venin equivalent.
\item \textbf{RC lowpass filter:} $R = 10\,\text{k}\Omega$, $C = 10\,$nF ($f_{3\text{dB}} \approx 1.6\,$kHz). Sweep 100\,Hz to 100\,kHz.
\item \textbf{Bridge rectifier with cap filter:} Vary cap and load, watch ripple change. Note narrow diode conduction angle.
\item \textbf{LC resonant tank:} $L = 100\,\mu$H, $C = 100\,$pF ($f_0 \approx 1.6\,$MHz). Sweep frequency, observe peak.
\end{enumerate}

\end{document}
